% Version 0.25. Complete exceptional odd-plus count supplier.
\section{Vanishing in the exceptional odd-plus case}
\label{sec:count-z-closure}

\status{CZ1--CZ6 prove the literal predicates \texttt{PropZ} and
\texttt{CountZ}, with the uniform exponent $C=0$.  The proof works on
the original algebraic closure and coefficient ring.  It proves the
two residual-character vanishing arguments, the needed cyclotomic
control, and the quadratic class-number estimate.  The accepted
arithmetic input is finite global class field theory in its ray-class
and Hilbert-class-field forms.  No local-duality, K-theory, automorphy,
Bernoulli-unit or other count predicate is retained as a premise.}

Fix an odd prime $\ell$, an algebraic closure $\Omega$ of
$\mathbb Q$, and a discrete valuation ring $O$ finite as a
$\mathbb Z_\ell$-algebra.  Put
\[
 G=\operatorname{Aut}_{\mathbb Q}(\Omega),\qquad
 M=O\otimes_{\mathbb Z_\ell}(\mathbb Q_\ell/\mathbb Z_\ell),
 \qquad \alpha=\chi_\ell^{-e}:G\longrightarrow O^\times.
 \tag{CZ.1}
\]
Here $M$ has the discrete topology, and $\chi_\ell$ is the
source's actual cyclotomic character, followed by the structure map
to $O$.  The action on $M$ is scalar multiplication by $\alpha$.
For a valuation $A$ of $\Omega$ above a finite rational prime,
write $D_A$ and $I_A$ for its decomposition and inertia groups.
Define the strict unramified Selmer subgroup
\[
 \begin{split}
 \mathcal S_e(O)={}&
 \ker\!\left(H^1(G,M(\alpha))\longrightarrow
     \prod_{A\mid\ell}H^1(D_A,M(\alpha))\right)
 \\
 &\cap\bigcap_A\ker\!\left(H^1(G,M(\alpha))
                       \longrightarrow H^1(I_A,M(\alpha))\right).
 \end{split}
 \tag{CZ.2}
\]
The source imposes its inertia condition for all its valuation
subrings; all finite-place conditions displayed here are among them.
We prove vanishing already under (CZ.2), which therefore proves
vanishing for the source group as well.

\paragraph{Classical input boundary.}
We use the finite global reciprocity and existence theorem: a finite
abelian extension is a quotient of the ray class group for a modulus
supported on its ramified places, and the ray class field of modulus
$1$ has Galois group the ordinary ideal class group, with real places
remaining real.  The rational ray class field of modulus
$\ell^r\infty$ is $\mathbb Q(\mu_{\ell^r})$.
These are the statements in Milne,
\href{https://www.jmilne.org/math/CourseNotes/CFT.pdf}{\emph{Class
Field Theory}}, v4.03, V, Theorems~3.5--3.6, Corollary~3.7,
and Examples~3.9--3.10.  We explain below the transfer from each
actual cocycle to the relevant finite extension.  Basic integral
bases, ideal factorization, finite Galois ramification, Eisenstein's
criterion and Hensel lifting are used in their usual elementary forms;
see Milne,
\href{https://www.jmilne.org/math/CourseNotes/ANT.pdf}{\emph{Algebraic
Number Theory}}, v3.08, \S\S2--3 and 6--7.  In particular the
class-number inequality below is proved, not retained as a leaf.

\begin{proposition}[CZ1: the actual coefficients and small cocycles]
\label{prop:cz-coefficients}
The module $M$ in (CZ.1) is $\ell$-divisible and every element is
annihilated by a power of $\ell$.  Its $\ell$-torsion $M[\ell]$
is a finite-dimensional $\mathbb F_\ell$-vector space.  Every class
in $H^1(G,M(\alpha))$ is annihilated by a power of $\ell$.
Every class annihilated by $\ell$ has a continuous cocycle
representative taking values in $M[\ell]$.

The source's \texttt{Chr.V}$(\ell,e,1)$ is the $O$-length of
its subgroup of (CZ.2), on exactly these coefficients and this
negative cyclotomic action.  For a representative $c$, its local
conditions say that there is a possibly different $m_H\in M$ for
each required local group $H$, with
\[
                  c(g)=(\alpha(g)-1)m_H\quad(g\in H).
                  \tag{CZ.3}
\]
Consequently it suffices to show that the $\ell$-torsion of
$\mathcal S_e(O)$ is zero.
\end{proposition}

\begin{proof}
For the actual source coefficient $O=c.O$, the structure map
$\mathbb Z_\ell\to O$ is injective because $O$ is a subalgebra
of $\texttt{PadicAlgCl}(\ell)$.  More generally it is injective
under the stated discrete valuation ring hypotheses: otherwise
its kernel is the nonzero prime ideal $(\ell)$; the finite algebra
$O$ would then be a finite integral domain, hence a field,
contrary to the discrete valuation ring hypothesis.  Since $O$ is
a domain, it is torsion-free and finite over the principal ideal
domain $\mathbb Z_\ell$, hence free of some finite rank $a$.
As a $\mathbb Z_\ell$-module, $M$ is therefore a direct sum of
$a$ copies of $\mathbb Q_\ell/\mathbb Z_\ell$.  This proves its
divisibility, primary torsion and the asserted description of
$M[\ell]$; no freeness over $O/\ell O$ is required.

The action on each $M[\ell^r]$ factors through the reduction of
the continuous cyclotomic character modulo $\ell^r$, so it is
continuous.  A continuous cocycle from the compact group $G$ to
the discrete module $M$ has finite image.  A single power of
$\ell$ kills all its values, and therefore also its cohomology
class.  If $\ell[c]=0$, there is $m\in M$ such that
$\ell c(g)=(\alpha(g)-1)m$.  Choose $b\in M$ with $\ell b=m$.
Then
\[
                c'(g)=c(g)-(\alpha(g)-1)b
                \tag{CZ.4}
\]
is a continuous cohomologous cocycle with $\ell c'=0$.
Subtracting this coboundary preserves every local condition (CZ.3).

In the source, \texttt{QO} is the tensor product in (CZ.1),
\texttt{rho}$(\ell,e,1)$ is scalar multiplication by
$\chi_\ell^{-e}$, and \texttt{ShaU} at the full Galois group is
the intersection of the restriction kernels for decomposition
groups above $\ell$ and all inertia groups.  Its cohomology uses
the cocycle equation
$c(gh)=c(g)+\alpha(g)c(h)$ and coboundaries as in (CZ.3), so
these identifications preserve both the original module and the
local restrictions.  Its length is exactly \texttt{Chr.V}.
Finally, if a nonzero primary-torsion group element has least
annihilating power $\ell^r$, multiplication by $\ell^{r-1}$
gives a nonzero element killed by $\ell$.  Thus the last
assertion follows without any finiteness assumption on the Selmer
group.
\end{proof}

\begin{proposition}[CZ2: cyclotomic inertia and rational control]
\label{prop:cz-cyclotomic-control}
Put $G_\infty=\ker(\chi_\ell:G\to\mathbb Z_\ell^\times)$.
For every $A\mid\ell$, the map
$I_A\to\mathbb Z_\ell^\times$ induced by $\chi_\ell$ is
surjective.  At primes different from $\ell$ the character is
trivial on inertia.  Every continuous additive homomorphism
$f:G\to M[\ell]$ unramified outside $\ell$ is zero on
$G_\infty$.

If a class in $\mathcal S_e(O)$ restricts to zero on $G_\infty$,
it is zero already on $G$.  This is the exact injectivity needed
here; it involves no change of algebraic closure or coefficient
ring.
\end{proposition}

\begin{proof}
For $r\ge1$, the polynomial
$\Phi_{\ell^r}(1+T)$ is Eisenstein at $\ell$: its constant
term is $\ell$, and modulo $\ell$ it is
$T^{\ell^{r-1}(\ell-1)}$.  Thus adjoining a primitive
$\ell^r$th root to $\mathbb Q_\ell$ gives a totally ramified
extension of degree $\varphi(\ell^r)$.  Globally the degree is
at most this number and at least this local degree, so the
cyclotomic Galois group is the full
$(\mathbb Z/\ell^r\mathbb Z)^\times$, and inertia at every
chosen $A\mid\ell$ maps onto it.  For a prescribed
$u\in\mathbb Z_\ell^\times$, its finite-level inverse images
in the compact group $I_A$ are nonempty nested closed sets.
Their intersection gives an element whose actual cyclotomic
character is $u$.  This supplies compatible infinite-level values,
not just unrelated finite-level surjections.

At a prime $q\ne\ell$, choose a finite residue extension containing
the $\ell^r$th roots of unity.  Their derivatives in
$X^{\ell^r}-1$ are units, so Hensel lifting places all these
roots in the unramified extension with that residue field.  Inertia
therefore acts trivially on every $\mu_{\ell^r}$, and hence on
$\chi_\ell$.  These finite ramification statements apply to the
decomposition and inertia groups of the original valuations of
$\Omega$.

The finite image of $f$ cuts out a finite abelian extension of
$\mathbb Q$ inside the original $\Omega$.  It is unramified
outside $\ell$ by hypothesis; complex conjugation has zero image
since its order divides $2$ and $M[\ell]$ has odd exponent.
Global class field theory consequently places this field in the
ray class field of some modulus $\ell^r\infty$, hence in
$\mathbb Q(\mu_{\ell^r})\subset\Omega$.
The subgroup $G_\infty$ fixes that field, proving $f|_{G_\infty}=0$.

For the last assertion, choose a representative $c$ for the given
class.  The action on $G_\infty$ is trivial, so every coboundary
on that subgroup is identically zero.  Trivial restriction of the
class therefore means $c|_{G_\infty}=0$.  The cocycle equation
shows that $c(g)$ depends only on $\chi_\ell(g)$.
Fix $A\mid\ell$ and choose $b\in M$ from its decomposition
condition (CZ.3).  For arbitrary $g\in G$, choose
$t\in I_A\subset D_A$ with $\chi_\ell(t)=\chi_\ell(g)$.
Then $c(g)=c(t)=(\alpha(t)-1)b=(\alpha(g)-1)b$.
This is one global coboundary, as required.
\end{proof}

\begin{proposition}[CZ3: the quadratic field and its class number]
\label{prop:cz-quadratic-class-number}
Suppose $\ell\ge5$ and $\ell\equiv1\pmod4$.  Inside the
original $\Omega$ there is a real quadratic field
$K=\mathbb Q(\sqrt\ell)$ of discriminant $\ell$ such that
\[
 G_K=\ker\bigl(\eta\circ\overline\chi_\ell\bigr),
 \qquad h_K<\ell,
 \tag{CZ.5}
\]
where $\eta$ is the quadratic character of
$\mathbb F_\ell^\times$ and $h_K$ is the ordinary ideal class
number.  Every continuous additive homomorphism
$G_K\to M[\ell]$ which is unramified at all finite places is
zero.
\end{proposition}

\begin{proof}
Choose a primitive $\ell$th root $\zeta\in\Omega$, extend
$\eta$ by $\eta(0)=0$, and set
$\tau=\sum_{a\in\mathbb F_\ell}\eta(a)\zeta^a$.
If $g\zeta=\zeta^u$, reindexing the sum gives
$g\tau=\eta(u)\tau$.  Let
$\tau'=\sum_b\eta(b)\zeta^{-b}=\eta(-1)\tau$.
The change of variables $a=tb$ for $b\ne0$ gives
\[
 \begin{split}
 \tau\tau'
 &=\sum_{t\in\mathbb F_\ell}\eta(t)
                  \sum_{b\ne0}\zeta^{(t-1)b}\\
 &=(\ell-1)-\sum_{t\ne1}\eta(t)=\ell.
 \end{split}
 \tag{CZ.6}
\]
Here the inner sum is $\ell-1$ for $t=1$ and $-1$ otherwise,
and $\sum_t\eta(t)=0$.  Since $\eta(-1)=1$, we have
$\tau^2=\ell$.  A rational square cannot have odd valuation at
the prime $\ell$, so $\tau\notin\mathbb Q$.  Thus
$K=\mathbb Q(\tau)\subset\Omega$ is quadratic and real, and
the transformation formula gives the asserted subgroup $G_K$.

The element $\omega=(1+\tau)/2$ is integral, satisfying
$X^2-X+(1-\ell)/4=0$.  The discriminant of the order
$\mathbb Z[\omega]$ is $\ell$.  The square of its index in
$\mathcal O_K$ divides this prime, by the change-of-integral-basis
formula for discriminants.  The index is therefore one, proving
$\operatorname{disc}(K)=\ell$.

Here is the class-number bound in full.  A nonzero integral ideal
$\mathfrak a$ embeds in $\mathbb R^2$ as a lattice of covolume
$\sqrt\ell\,N\mathfrak a$, by the discriminant formula and
the index $[\mathcal O_K:\mathfrak a]=N\mathfrak a$.
The symmetric diamond $|x_1|+|x_2|<t$ has area $2t^2$.
If this exceeds four times the covolume, it contains a nonzero
lattice point: half the diamond has area greater than the
covolume, so two of its points have the same image in a
fundamental parallelogram, and their difference is a nonzero
lattice point in the diamond.  This is the required
two-dimensional convex-body argument.  Consequently there is
$0\ne a\in\mathfrak a$ with
\[
 |N_{K/\mathbb Q}a|\le\tfrac12\sqrt\ell\,N\mathfrak a.
 \tag{CZ.7}
\]
Indeed $|x_1x_2|\le t^2/4$, and the bound at the limiting
value follows because nonzero absolute norms are positive integers.
The integral ideal $(a)\mathfrak a^{-1}$ represents the inverse
class and has norm at most $\sqrt\ell/2$.  Every ideal class
therefore has an integral representative of norm at most
$X=\sqrt\ell$.

For an integer $n\ge1$, the number of ideals of norm $n$ in a
quadratic field is at most the divisor number $d(n)$.
At a rational prime $q$, the split case has $r+1$ ideals of norm
$q^r$, the ramified case has one, and the inert case has at most
one.  Unique ideal factorization multiplies these bounds.
It follows that
\[
 \begin{split}
 h_K&\le\sum_{n\le X}d(n)
     =\sum_{a\le X}\lfloor X/a\rfloor
     \le X\sum_{a\le X}\frac1a\\
    &\le X(1+\log X)<X^2=\ell.
 \end{split}
 \tag{CZ.8}
\]
The harmonic-sum inequality follows by comparison with the integral
of $1/t$, and the strict inequality uses $\log X<X-1$ for
$X>1$.

Finally, the finite image of an unramified homomorphism as stated
defines a finite abelian extension of $K$.  At real places its
inertia has order at most two and therefore maps to zero in
$M[\ell]$.  The extension is consequently a quotient of the
ordinary Hilbert class field.  Its image is annihilated by $h_K$
as well as by $\ell$.  Since $0<h_K<\ell$, these two integers
are coprime, and the image is zero.
\end{proof}

\begin{proposition}[CZ4: residual-trivial vanishing]
\label{prop:cz-trivial-vanishing}
For $e\ge1$ with $\ell-1\mid e$, one has
$\mathcal S_e(O)=0$ and hence
$\texttt{Chr.V}(\ell,e,1)=0$.
\end{proposition}

\begin{proof}
Let $x\in\mathcal S_e(O)$ be annihilated by $\ell$ and use
CZ1 to represent it by an $M[\ell]$-valued cocycle $c$.
For every $g\in G$, the reduction of $\chi_\ell(g)^{-e}$
modulo $\ell$ is one because $\ell-1\mid e$.
Thus $\alpha(g)-1$ is a multiple of $\ell$ in $O$ and the
action on $M[\ell]$ is trivial.  The cocycle $c$ is therefore
an additive homomorphism.
At every $q\ne\ell$, CZ2 says $\alpha$ is trivial on
$I_A$; its local coboundary condition (CZ.3) then says
$c|_{I_A}=0$.  It is an unramified-outside-$\ell$ homomorphism,
so CZ2 gives $c|_{G_\infty}=0$.  The injectivity proved in
the same proposition gives $x=0$.  By CZ1 this kills the whole
Selmer group.  Its original $O$-length is consequently zero.
\end{proof}

\begin{proposition}[CZ5: residual-quadratic vanishing]
\label{prop:cz-quadratic-vanishing}
Suppose $\ell\ge5$, $e$ is even,
$\ell-1\mid2e$ and $\ell-1\nmid e$.  Then
$\mathcal S_e(O)=0$ and hence
$\texttt{Chr.V}(\ell,e,1)=0$.
\end{proposition}

\begin{proof}
Write $\ell-1=2a$.  The divisibility assumptions give
$e=ak$ with $k$ odd.  Since $e$ is even, $a$ is even and
$\ell\equiv1\pmod4$.  The residual character of
$\alpha=\chi_\ell^{-e}$ is exactly
$\eta\circ\overline\chi_\ell$, so it is trivial on the
actual $G_K$ in CZ3.  In particular $G_K$ acts trivially on
$M[\ell]$.

Choose $u\in\mathbb Z_\ell^\times$ with nonsquare reduction.
For any $g$ with $\chi_\ell(g)=u$, the scalar
$\alpha(g)-1$ is a unit of $O$: it is the image of the unit
$u^{-e}-1\in\mathbb Z_\ell$, whose reduction is $-2\ne0$.
Fix one place $A_0\mid\ell$ and, using CZ2, choose
$\delta\in I_{A_0}$ with cyclotomic value $u$.  Then
$\delta\notin G_K$.

Let $x\in\mathcal S_e(O)$ be killed by $\ell$ and choose an
$M[\ell]$-valued representative $c$ by CZ1.  Subtract the
coboundary of
$b=(\alpha(\delta)-1)^{-1}c(\delta)\in M[\ell]$.
The resulting representative, still denoted $c$, has
\[
                  c(\delta)=0,\qquad \ell c(g)=0
                  \quad(g\in G).
                  \tag{CZ.9}
\]
It retains all its local conditions.

Fix any $A\mid\ell$ and its local witness $m_A$ from (CZ.3)
on $D_A$.  Choose $t\in I_A$ with $\chi_\ell(t)=u$.
Then
\[
       0=\ell c(t)=(\alpha(t)-1)\ell m_A
       \quad\Longrightarrow\quad \ell m_A=0,
       \tag{CZ.10}
\]
because $\alpha(t)-1$ is a unit.  For $g\in D_A\cap G_K$
the action on $m_A\in M[\ell]$ is trivial; hence
$c(g)=0$.  At primes different from $\ell$, $\alpha$ is
trivial on inertia and (CZ.3) likewise gives $c|_{I_A}=0$.
Restriction of $c$ to $G_K$ is an additive homomorphism, since
the coefficient action is trivial there.  The preceding statements
make it unramified at every finite place of $K$ (and even zero
on its decomposition groups above $\ell$).  CZ3 therefore gives
$c|_{G_K}=0$.

The subgroup $G_K$ has index two and $\delta$ lies in the other
coset.  For any $g\notin G_K$ write $g=\delta h$ with
$h\in G_K$; then
$c(g)=c(\delta)+\alpha(\delta)c(h)=0$ by (CZ.9).
Thus $c=0$, so $x=0$.  Applying the primary-torsion argument
of CZ1 proves the full vanishing and the stated length equality.
\end{proof}

\begin{theorem}[CZ6: the literal exceptional count, uniformly with exponent zero]
\label{thm:cz-count-z}
The source predicates \texttt{PropZ} and \texttt{CountZ} hold.
In \texttt{CountZ}, one may choose $C=0$ for every primitive
character, before choosing the host, twist cutoff, algebraic
closure, prime, coefficient ring, twist or character in the host.
In fact the counted length vanishes identically in this case.
\end{theorem}

\begin{proof}
For \texttt{PropZ}, fix its original $\Omega,\ell,c,e$ with
$\ell\ge5$, $e\ge2$, $e$ even and $\ell-1\mid2e$;
here $O=c.O$.  If $\ell-1\mid e$, apply CZ4.  Otherwise
apply CZ5.  Both return the exact length
$\texttt{Chr.V}(\ell,e,1)=0$ on the original objects.

Now fix any data in the universal statement \texttt{CountShape}
for the case \texttt{CaseZ}.  By its definition this case says
\[
          \ell\notin B_\chi,\qquad
          \ell-1\mid2e,\qquad \psi=1.
          \tag{CZ.11}
\]
The exceptional prime set $B_\chi$ contains $3$; the count
statement already assumes the prime $\ell\ne2$.  Thus
$\ell\ge5$.  Membership $\psi\in h.X(c,e)$ includes
$\psi(-1)=(-1)^e$.  Because $\psi=1$ and $O$ has
characteristic zero, this forces $e$ even; the given $e>0$
then gives $e\ge2$.  The first part of the theorem applies.

The Galois character attached to the trivial Dirichlet character
is again trivial.  The source's \texttt{Vlen} in this case is
$\texttt{Chr.V'}(\ell,h.N,e,1)$, whose local conditions impose
decomposition restrictions above $\ell h.N$.  These include
all decomposition restrictions above $\ell$ and the same
inertia restrictions, so its Selmer group is a subgroup of the
zero group just proved.  Hence
\[
 \texttt{Vlen}(\Omega,c,e,1)=0,
 \qquad \texttt{Vn}(\Omega,c,e,1)=0.
 \tag{CZ.12}
\]
The first equality proves finiteness; the second follows from
the source's normalization of the length, whose numerator is
zero.  Therefore its actual predicate \texttt{VLe} holds with
left side $\texttt{Vn}\log\ell=0$ and right side
$((\ell:\mathbb R)h.N e)^0=1$.
Choosing $C=0$ before all the universally quantified data proves
\texttt{CountZ}.  No threshold, host, splitting field, coefficient
extension or discriminant is replaced in this last step.
\end{proof}

\paragraph{Closed dependency and remaining work.}
This closes one entire input of the source's implication from the
counts to \texttt{PointHV}.  The other count suppliers
\texttt{CountI}, \texttt{CountII}, \texttt{CountIII} and
\texttt{CountV} retain their own obligations; in particular this
chapter does not supply the automorphy and host-rigidity inputs
used for \texttt{CountV}.  The current vanishing proof is independent
of those inputs and of \texttt{LemmaK}.
